\documentclass[10pt,a4paper]{article}

% Layout and language
\usepackage[margin=1in]{geometry}
\usepackage[british]{babel}
\usepackage[T1]{fontenc}
\usepackage[utf8]{inputenc}
\usepackage{microtype}

% Spacing and front matter
\usepackage{setspace}
\usepackage{authblk}
\usepackage{lineno}

% Mathematics
\usepackage{mathtools,amssymb}
\usepackage{siunitx}

% Tables and figures
\usepackage{graphicx}
\usepackage{booktabs}
\usepackage{tabularx}
\usepackage{adjustbox}
\usepackage{caption}
\usepackage{subcaption}
\usepackage{float}

% References and links
\usepackage{natbib}
\usepackage[hidelinks]{hyperref}
\usepackage[nameinlink,noabbrev]{cleveref}

% Sensible defaults
\sisetup{detect-all}
\onehalfspacing
% \linenumbers

\title{Derivation of the epilink compatibility score}
\author{}
\date{}

\begin{document}

\maketitle

For a sampled pair $(i,j)$, let $t_{ij}$ denote the sampling-time difference and $g_{ij}$ the observed consensus-level genetic distance. The \texttt{epilink} framework compares a finite set of recent latent transmission scenarios linking the pair, computes a compatibility score for each scenario, and reports the compatibility of a user-specified target scenario.

\subsection*{Observed quantities and timing model}

Let $t_{\mathrm{samp},i}$ and $t_{\mathrm{samp},j}$ denote the sampling dates, measured in days on a common time scale, and define
\begin{align}
t_{ij} = t_{\mathrm{samp},j} - t_{\mathrm{samp},i}.
\end{align}
Let $g_{ij}$ denote the observed consensus-level genetic distance, measured in nucleotide differences.

The natural-history model follows the variable infectiousness $E/P/I$ framework of \citet{Hart2021}, in which infection is partitioned into latent ($E$), presymptomatic infectious ($P$), and symptomatic infectious ($I$) stages with durations $y_E$, $y_P$, and $y_I$. We additionally introduce a sampling delay $y_S$, defined as the time from symptom onset to sample collection, and write $t_{\mathrm{toit}}$ for the time from infectiousness onset to transmission. The incubation period and generation time are
\begin{align}
\tau_{\mathrm{inc}} &= y_E + y_P,\\
\tau_{\mathrm{gen}} &= y_E + t_{\mathrm{toit}},
\end{align}
so that the total time from infection to sampling is $\tau_{\mathrm{inc}} + y_S$.

Following~\cite{Hart2021}, the stage durations are modelled as independent Gamma random variables, with $E$ and $P$ sharing a common scale:
\begin{align}
y_E &\sim \mathrm{Gamma}(k_E,\theta_{\mathrm{inc}}),\\
y_P &\sim \mathrm{Gamma}(k_P,\theta_{\mathrm{inc}}),\\
y_I &\sim \mathrm{Gamma}(k_I,\theta_I),
\end{align}
where
\begin{align}
k_P = k_{\mathrm{inc}} - k_E,
\qquad
\tau_{\mathrm{inc}} \sim \mathrm{Gamma}(k_{\mathrm{inc}},\theta_{\mathrm{inc}}).
\end{align}
Using the parameterisation of Hart et al.,
\begin{align}
k_{\mathrm{inc}} &= 5.807, &
\theta_{\mathrm{inc}} &= 0.948, &
k_I &= 1,\\
k_E &= 3.38, &
\mu &= 0.37, &
\alpha &= \frac{\beta_P}{\beta_I} = 2.29,
\end{align}
with $\theta_I = 1/\mu$.

Transmission timing is induced by the infectiousness profile rather than specified independently. Writing $f_P$ and $F_P$ for the density and cumulative distribution functions of the presymptomatic period, the induced density of $t_{\mathrm{toit}}$ is
\begin{equation}
\label{eq:toit_density}
f_{\mathrm{toit}}(t_{\mathrm{toit}})
=
C\!\left[
\alpha \bigl(1-F_P(t_{\mathrm{toit}})\bigr)
+
\int_{0}^{t_{\mathrm{toit}}}
\bigl(1-F_P(t_{\mathrm{toit}}-y_P)\bigr)f_P(y_P)\,dy_P
\right],
\end{equation}
where $C$ is a normalising constant. To connect this model to observed sampling dates, \texttt{epilink} adds a sampling delay
\begin{align}
y_S \sim \mathrm{Gamma}(k_S,\theta_S).
\end{align}
All timing quantities are treated as latent random variables and integrated over by Monte Carlo.

\subsection*{Latent transmission scenarios}

Let $M$ denote the maximum hidden depth allowed in the latent relationship between cases $i$ and $j$. The candidate set is
\begin{align}
\mathcal{S}_M
=
\left\{
H_{\mathrm{AD}}(m): m=0,\ldots,M
\right\}
\cup
\left\{
H_{\mathrm{CA}}(m_i,m_j): m_i,m_j \ge 0,\ m_i+m_j \le M
\right\}.
\end{align}
Here, $H_{\mathrm{AD}}(m)$ denotes an ancestor--descendant scenario in which case $j$ descends from case $i$ through $m$ unsampled intermediates, whereas $H_{\mathrm{CA}}(m_i,m_j)$ denotes a common-ancestor scenario in which both sampled cases descend from an unsampled source with branch depths $m_i$ and $m_j$. The cases $H_{\mathrm{AD}}(0)$ and $H_{\mathrm{CA}}(0,0)$ correspond to direct transmission and co-primary infection, respectively.

\subsection*{Temporal draws}

For a given scenario $s \in \mathcal{S}_M$, let $t_x$ denote the infection time of the latent reference point: the infection time of case $i$ under an ancestor--descendant scenario, or the infection time of the shared unsampled source under a common-ancestor scenario. Let $A_k(s)$ denote the elapsed time from this reference point to infection of sampled case $k \in \{i,j\}$. Then
\begin{align}
t_{\mathrm{samp},k}
=
t_x + A_k(s) + \tau_{\mathrm{inc},k} + y_{S,k},
\end{align}
so the model-implied sampling-time difference is
\begin{align}
T_s
=
A_j(s) - A_i(s)
+
\tau_{\mathrm{inc},j} + y_{S,j}
-
\tau_{\mathrm{inc},i} - y_{S,i}.
\label{eq:tij_master}
\end{align}

Under $H_{\mathrm{AD}}(m)$,
\begin{align}
A_i\!\left(H_{\mathrm{AD}}(m)\right) &= 0,\\
A_j\!\left(H_{\mathrm{AD}}(m)\right) &= \sum_{r=0}^{m}\tau_{\mathrm{gen},r},
\end{align}
and therefore
\begin{align}
T_{ij}\!\left(H_{\mathrm{AD}}(m)\right)
=
\sum_{r=0}^{m}\tau_{\mathrm{gen},r}
+
\tau_{\mathrm{inc},j} + y_{S,j}
-
\tau_{\mathrm{inc},i} - y_{S,i}.
\label{eq:T_AD}
\end{align}

Under $H_{\mathrm{CA}}(m_i,m_j)$,
\begin{align}
A_i\!\left(H_{\mathrm{CA}}(m_i,m_j)\right)
&=
\sum_{r=1}^{m_i+1}\tau_{\mathrm{gen},r}^{(i)},\\
A_j\!\left(H_{\mathrm{CA}}(m_i,m_j)\right)
&=
\sum_{s=1}^{m_j+1}\tau_{\mathrm{gen},s}^{(j)},
\end{align}
so that
\begin{align}
T_{ij}\!\left(H_{\mathrm{CA}}(m_i,m_j)\right)
=
\sum_{s=1}^{m_j+1}\tau_{\mathrm{gen},s}^{(j)}
-
\sum_{r=1}^{m_i+1}\tau_{\mathrm{gen},r}^{(i)}
+
\tau_{\mathrm{inc},j} + y_{S,j}
-
\tau_{\mathrm{inc},i} - y_{S,i}.
\label{eq:T_CA}
\end{align}
For each scenario $s$, Monte Carlo simulation yields draws $T_s^{(1)},\ldots,T_s^{(N)}$ from the induced temporal distribution.

\subsection*{Genetic draws}

The genetic component assumes that within-host divergence is negligible on the timescales of interest and that consensus divergence accumulates along transmission-linked lineages. Under scenario $s$, let $B_s$ denote the effective transmission-related branch length, in days, separating the sampled genomes.

For ancestor--descendant scenarios,
\begin{align}
B_{H_{\mathrm{AD}}(m)}
=
\sum_{r=0}^{m}\tau_{\mathrm{gen},r}
+
\tau_{\mathrm{inc},j} + y_{S,j}
-
\left(\tau_{\mathrm{inc},i} + y_{S,i}\right).
\label{eq:B_AD}
\end{align}
For common-ancestor scenarios,
\begin{align}
B_{H_{\mathrm{CA}}(m_i,m_j)}
=
\sum_{r=1}^{m_i+1}\tau_{\mathrm{gen},r}^{(i)}
+
\sum_{s=1}^{m_j+1}\tau_{\mathrm{gen},s}^{(j)}
+
\tau_{\mathrm{inc},i} + y_{S,i}
+
\tau_{\mathrm{inc},j} + y_{S,j}.
\label{eq:B_CA}
\end{align}
For each scenario $s$, Monte Carlo simulation yields draws $B_s^{(1)},\ldots,B_s^{(N)}$.

Let $r$ denote the median substitution rate per site per year and $L$ the genome length in sites. The corresponding per-genome daily substitution rate is
\begin{align}
\lambda = \frac{rL}{365}.
\label{eq:lambda_strict}
\end{align}
Under a relaxed uncorrelated log-normal clock (UCLN), one instead samples
\begin{align}
r^{(n)} \sim \mathrm{LogNormal}(\mu_r,\sigma_r^2),
\end{align}
with $\mu_r$ chosen so that the median equals the specified substitution rate $r$, and sets
\begin{align}
\lambda^{(n)} = \frac{r^{(n)}L}{365}.
\label{eq:lambda_relaxed}
\end{align}

Given branch-length draw $B_s^{(n)}$ and clock-rate draw $\lambda^{(n)}$, the genetic draw used for scoring is
\begin{align}
G_s^{(n)} =
\lambda^{(n)} B_s^{(n)}
\end{align}
under a deterministic mutational process, and
\begin{align}
G_s^{(n)} \mid B_s^{(n)}, \lambda^{(n)}
\sim
\mathrm{Poisson}\!\left(\lambda^{(n)} B_s^{(n)}\right)
\end{align}
under a stochastic mutational process.

\subsection*{Compatibility scores}

For any observed quantity $x_{\mathrm{obs}}$ and scenario $s$ with Monte Carlo draws $x_1^{(s)},\ldots,x_{n_s}^{(s)}$, define the percentile score
\begin{align}
p_s(x_{\mathrm{obs}})
=
\frac{\#\left\{i : x_i^{(s)} \le x_{\mathrm{obs}}\right\}}{n_s},
\end{align}
and the corresponding compatibility score
\begin{align}
C_s(x_{\mathrm{obs}})
=
1 - 2\left|p_s(x_{\mathrm{obs}}) - 0.5\right|.
\end{align}
Applying this definition to the temporal draws gives
\begin{align}
p_{T,s}(i,j)
=
\frac{\#\left\{n : T_s^{(n)} \le t_{ij}\right\}}{N},
\qquad
C_{T,s}(i,j)
=
1 - 2\left|p_{T,s}(i,j) - 0.5\right|,
\end{align}
and applying it to the genetic draws gives
\begin{align}
p_{G,s}(i,j)
=
\frac{\#\left\{n : G_s^{(n)} \le g_{ij}\right\}}{N},
\qquad
C_{G,s}(i,j)
=
1 - 2\left|p_{G,s}(i,j) - 0.5\right|.
\end{align}
The overall compatibility assigned to scenario $s$ is the product of the temporal and genetic compatibilities,
\begin{align}
C_s(i,j) = C_{T,s}(i,j)\,C_{G,s}(i,j).
\label{eq:scenario_compatibility}
\end{align}

\subsection*{Target subset and computation}

Let $\mathcal{S}_\star \subseteq \mathcal{S}_M$ denote a user-specified target subset chosen when the \texttt{epilink} object is constructed. The reported target score is
\begin{align}
\mathrm{score}_{\mathcal{S}_\star}(i,j) = \sum_{s \in \mathcal{S}_\star} C_s(i,j).
\end{align}
When $\mathcal{S}_\star$ contains a single element, this reduces to the original single-scenario score. The implementation also returns $C_s(i,j)$ for every $s \in \mathcal{S}_M$.

For each scenario $s \in \mathcal{S}_M$, \texttt{epilink} precomputes and stores $N$ draws of $T_s$ and $B_s$ at construction time. It then converts the cached branch-length draws into cached mutation draws $G_s^{(n)}$ using the chosen mutational process. Each pairwise evaluation therefore only requires percentile calculations against the cached draws, which is well suited to large datasets.

\bibliographystyle{plainnat}
\bibliography{references}

\end{document}
